% concatenated from fujita_on_general_hk_Alessandro_Pilastro.tex
\documentclass{article}
\usepackage{graphicx} % Required for inserting images
\usepackage{amsmath}
\usepackage{amssymb}
\usepackage{amsopn}
\usepackage{amsthm}
\usepackage{amsfonts}
\usepackage[right]{eurosym}
\usepackage{mathtools}
\usepackage{comment}
\usepackage{mathrsfs}
\usepackage{hyperref}

\usepackage{color} 
\usepackage{graphicx}

\usepackage{tikz-cd}

%Letters
\DeclareMathOperator{\bC}{\mathbb{C}}
\DeclareMathOperator{\bZ}{\mathbb{Z}}
\DeclareMathOperator{\bN}{\mathbb{N}}
\DeclareMathOperator{\bP}{\mathbb{P}}
\DeclareMathOperator{\cA}{\mathcal{A}}
\DeclareMathOperator{\cF}{\mathcal{F}}
\DeclareMathOperator{\cB}{\mathcal{B}}
\DeclareMathOperator{\cP}{\mathcal{P}}
\DeclareMathOperator{\cV}{\mathcal{V}}
\DeclareMathOperator{\cO}{\mathcal{O}}

%New math operator
\DeclareMathOperator{\Hilb}{Hilb}
\DeclareMathOperator{\Sym}{Sym}
\DeclareMathOperator{\Pic}{Pic}
\DeclareMathOperator{\Kum}{Kum}
\DeclareMathOperator{\Div}{div}
\DeclareMathOperator{\gr}{\mathbf{Gr}}
\DeclareMathOperator{\Id}{Id}
\newcommand{\hilb}[2][n]{#2^{[#1]}}
\newcommand{\sym}[2][n]{#2^{(#1)}}

%Theorems

\newtheorem{proposition}{Proposition}[section]
\newtheorem{remark}[proposition]{Remark}
\newtheorem{lemma}[proposition]{Lemma}
\newtheorem{corollary}[proposition]{Corollary}
\newtheorem{theorem}[proposition]{Theorem}

\title{On Fujita's conjecture for a general hyperk\"ahler manifold in the standard series of examples}
\author{Alessandro Pilastro}
\date{}



\begin{document}

\maketitle
\begin{abstract}
   For the moduli spaces $\Sigma_{d,t}^n$ and $\Upsilon_{d,t}^n$ of polarized hyperk\"ahler manifolds of $\hilb{K3}$-type and $\Kum^n$-type respectively, with polarization with square $2d$ and divisibility $t$, we study general base point freeness and very ampleness of the polarization. We provide cases where these moduli space are connected and a formula characterizing when these spaces are non-empty.
\end{abstract}

\section*{Introduction}

Fujita's conjecture states that given an ample holomorphic line bundle $L$ on a compact complex manifold $X$ of complex dimension $n$, the line bundle $K_X\otimes L^{\otimes m}$ is base point free for all $m\geq n+1$ and very ample for $m\geq n+2$. In general, the conjecture is widely open, as only in low dimensions the conjecture is known to hold and in some specific examples of manifolds. Two important families that satisfy the conjecture are $K3$ surfaces and Abelian varieties, where the conjecture holds with lower bounds of $2$ and $3$, respectively. These values are stronger compared to what the conjecture suggests, so it would be interesting to study the conjecture on hyperk\"ahler manifolds.

Let $\Sigma_{d,t}^n$ be the moduli space polarized of hyperk\"ahler manifolds of $\hilb{K3}$-type with polarization of square $2d$ and divisibility $t$; let $\Upsilon_{d,t}^n$ be the analogous moduli space for $\Kum^n$-type manifolds. In \cite{riess}, U.Rie{\ss} proves that the general element of $\Sigma_{d,1}^2$ and $\Sigma_{d,2}^2$ is base point free, while in \cite{Varesco_2023}, M.Varesco proves a similar statement for the $\Kum^n$-type for $n=2,3,4$ and all possible divisibilities, except for a finite number of cases. The results of U.Rie{\ss} is generalized by O.Debarre in \cite{debarre2020hyperkahlermanifolds}, where bounds in terms of $d$ and $n$ are given for the general element of the moduli spaces $\Sigma_{d,1}^n$ and $\Sigma_{d,2}^n$ to be base point free and very ample. 

Inspired by these results, we formalize the techniques used in the aforementioned works to generalize the results in higher dimensions for all possible divisibility in both the case of $\hilb{K3}$-type and $\Kum^n$-type hyperk\"ahler. 

\begin{theorem}\label{theo_main}
    Let $n,t,d$ be positive integers with $n\geq 2$ and let $\tau:=\frac{t^2}{2(t-1)}$. Then the following holds:
    \begin{itemize}
        \item let $(X,H)$ be a general element of $\Sigma_{d,n}^1$ or $\Upsilon_{d,n}^1$ then $H$ is base point free and $H^{n+2}$ is very ample. Furthermore, if $d\geq n+1$ then the polarization of a general element of $\Sigma_{d,n}^1$ is very ample, if $d\ge n+4$ then the polarization of a general element of $\Upsilon_{d,n}^1$ is very ample;
        \item if $d\geq (\tau-1)n+\tau+1$ there exists a connected component of $\Sigma_{d,n}^t$ such that the polarization of a general element $(X,H)$ is base point free and $H^{n+2}$ is very ample;
        \item if $d\geq (\tau-1)n+2\tau+1$ there exists a connected component of $\Sigma_{d,n}^t$ such that the polarization of a general element is very ample;
        \item if $d\geq (\tau-1)n+2\tau-1$ there exists a connected component of $\Upsilon_{d,n}^t$ such that the polarization of a general element $(X,H)$ is base point free and $H^{n+2}$ is very ample;
        \item if $d\geq (\tau-1)n+3\tau-1$ there exists a connected component of $\Upsilon_{d,n}^t$ such that the polarization of a general element is very ample;
        These hold unless the considered moduli space is empty.
    \end{itemize}
\end{theorem}

In more interesting situations, to which we can apply the theorem, occur when these moduli spaces have one connected component. In \cite{onorati2022connectedcomponentsmodulispaces} and \cite{zbMATH06387271}, the authors provide a formula to compute the numbers of connected components of these moduli spaces. For example, with divisibilities $t=1,2$ and $t=p^a$, with $p>2$ prime, and $p^{a+1}$ not dividing $n$ the moduli spaces $\Sigma_{d,t}^n$ and $\Upsilon_{d,t}^n$ are connected. We conclude by providing, in Proposition \ref{prop_char_of_non_emptyness}, a characterization of when the considered moduli spaces are non-empty. This allows us to improve the inequalities presented in Theorem \ref{theo_main} in any specific case. 

The structure of the article is the following.

In Section \ref{sect_tautological}, we recall the main properties of the Hilbert scheme of points on a surface, with a focus on the tautological vector bundle $\hilb{L}$ on $\hilb{S}$ associated to any line bundle $L$ on the surface $S$. We prove how this construction preserves base point freeness. 

In Section \ref{sect_grass}, we provide the general notion of embeddings into Grassmannians via suitable vector bundles. In particular, we show how the tautological vector bundle $\hilb{L}$ defines a map into a Grassmannian if $L$ is $(n-1)$-very ample, which is an embedding if $L$ is $n$-very ample. Taking the determinant of $\hilb{L}$ allows us to obtain a base point free and a very ample line bundle on $\hilb{S}$.

Sections \ref{sect_hyper} and \ref{sect_moduli} recall the definition of hyperk\"ahler manifolds and provide the statements of some properties that we will use throughout the article. In particular, we present the lattice structure of the Picard group of $\hilb{K3}$-type and $\Kum^n$-type hyperk\"ahler manifolds. We also show why base point freeness and very ampleness are open properties in families for hyperk\"ahler manifolds.

In section \ref{sect_main_theorems}, we combine the results to obtain the main theorem. The main idea of the proof is to build an element of the moduli space with the required properties. The openness of base point freeness and the very ampleness allows us to conclude that the general elements in the same connected component share the same properties. 

\section{Tautological vector bundles on $\Hilb^n$ of a surface}\label{sect_tautological}

Let $S$ be a smooth, projective complex surface and let $\hilb{S}:=\Hilb^n(S)$ the Hilbert scheme parameterizing $0-$dimensional subschemes of length $k$ of $S$. For an introduction on the subject one can refer, among other sources to \cite{nakajima_hilbert_schemes}.  Since $\hilb{S}$ is a fine moduli space we have a universal family $$\Xi_k=\{(x,\xi)\in S\times \hilb{S}|x\in \xi\}\subset S\times \hilb{S}$$ together with the natural projections $$S\xleftarrow{\pi_S}\Xi_k\xrightarrow{\pi_{\hilb{S}}}\hilb{S}.$$
Given any coherent sheaf on $F$ on $S$ we can construct the so called \emph{tautological sheaf associated to $F$} on $\hilb{S}$ as follows: $$\hilb{F}:=\pi_{\hilb{S}*}(\pi^*_S(F)).$$
Since the projection $\pi_{\hilb{S}}$ is flat of degree $k$, if $F$ is a vector bundle of rank $m$,  then the associated tautological bundle $\hilb{F}$ is a vector bundle of rank $mk$ with fibers
\begin{align}\label{eq_fiber_of_F^[n]}
    \hilb{F}(\xi)\cong H^0(\xi,F|_\xi),\ \forall \xi\in\hilb{S}.
\end{align}

The tautological bundles have the following properties.

\begin{proposition}[\cite{bini2022familiesbigstablebundles}, (4.5)]\label{prop_zeroth_cohom_F^[n]}
Let $S$ be a smooth, irreducible projective surface and let $F$ be a coherent sheaf on $S$ then $H^0(\hilb{S},\hilb{F})\cong H^0(S,F)$.
\end{proposition}

Recall that $F$ is $(n-1)$-\emph{very ample} on $S$ if the natural evaluation map $$H^0(S,F)\to H^0(S,F\otimes \cO_\xi)$$ is surjective for all $0$-dimensional subschemes $\xi$ of length $n$.

\begin{proposition}[\cite{bini2022familiesbigstablebundles}, Prop 4.5]\label{prop_tautological_globally_generated}
Let $S$ be a smooth, irreducible projective surface and let $F$ be a rank $r$ vector bundle on $S$ with $r\geq 1$. Let $n\geq 2$ be an integer. Then $F$ is $(n-1)$-very ample on $S$ if and only if the tautological bundle $\hilb{F}$ is globally generated.
\end{proposition}

Denoting the $k$-th cartesian product $S^k=S\times\dots\times S$ and the $k$-th symmetric product $\sym{S}=S^n/\Sym(n)$, we have the \emph{Hilbert-Chow morphism} $$\mu:\hilb{S}\to\sym{S}$$ given by $\mu(\xi)=\sum_{x\in\text{Supp($\xi$)}}\text{ length}_\xi(x)x$.
Let $L\in \Pic(S)$, denoting by $\pi_i:S^n\to S$ the $i$-th natural projection, we can define the line bundle $L^{\boxtimes n}:=\otimes_{i=1}^n\pi^*_i(L)$ which descends to a line bundle $\sym{L}\in\Pic(\sym{S})$. This provides a morphism $D_n:\Pic(S)\to \Pic(\hilb{S})$ by $D_n(L)=L_n$ where $L_n:=\mu^*(\sym{L})$.    

Using this morphism, we can explicitly write what the determinant line bundle of a tautological vector bundle is. Let $E$ be the exceptional divisor of the Hilbert-Chow morphism, and let $\delta\in  \Pic(\hilb{S})$ such that $2\delta=E$.

\begin{proposition}[\cite{Lehn_1999}, Lemma 3.7 - Thm 4.6]\label{prop_det_of_tautological}
    Let $S$ be a smooth surface and let L be a line bundle on $S$, then the line bundle $\det \hilb{L}=L_n-\delta$ on $\hilb{S}$.
\end{proposition}

To conclude this section, let us show that the morphism $D_n$ preserves base point freeness.

\begin{lemma}\label{lemma_bpf_of_L_n_on_Hilb_n}
Let $S$ be a smooth projective complex surface, and let $L\in\Pic(S)$ be base point free, then $L_n\in \Pic(\hilb{S})$ is also base point free.    
\end{lemma}
\begin{proof}
Let $\xi\in\hilb{S}$ and let $x_1,\dots,x_n$ be the points in the support of $\xi$ with multiplicities. We can find a section $s\in H^0(S,L)$ such that $s(x_i)\neq 0$ for all $i=1,\dots,n.$ This follows by induction on $n$. The base case is the definition of base point free line bundle. From the induction hypothesis, we have a section $\hat{s}$ such that $\hat{s}(x_i)\neq 0$ for $i=1,\dots,n-1$. If $\hat{s}(x_n)\neq 0$ then we are done. Suppose then that $\hat{s}(x_n)= 0$ and given $\tilde{s}$ a section non-vanishing at the point $x_n$, define $s=\alpha \hat{s}+\tilde{s}$. For a general $\alpha\in\bC$ the section $s$ satisfies the required property.

So, given $\xi\in\hilb{S}$ let $s$ be a section of $L$ non-vanishing at the points in the support of $\xi$. Consider now the section $\otimes_{i=1}^n\pi_i^*(s)\in H^0(S^n,L^{\boxtimes n})$ which is invariant under the action of $\Sym(n)$. It therefore descends to a section $\sym{s}\in H^0(\sym{S},\sym{L})$. Then $\sym{s}$ takes the value $\prod_{i=1}^ns(x_i)\neq0$ at the point $\sum_{i=1}^nx_i\in\sym{S}$. The pullback section $\mu^*(\sym{s})=\hilb{s}\in H^0(\hilb{S},L_n)$ does not vanish at the point $\xi\in\hilb{S}$ as required.

\end{proof}

 \section{Embeddings of Hilbert schemes into Grassmannians}\label{sect_grass}

We start by setting some notations for the Grassmannians. Let $V$ be a complex $d$-dimensional vector space, with $\gr(V,n)$ we denote the space of $n$-dimensional subspaces of $V$. By $U$ we denote the universal subbundle and by $Q$ the universal quotient bundle, which naturally fits in the following short exact sequence $$0\to U\to \gr(V,n)\times V\to Q\to 0.$$ 

Let $X$ be a compact complex manifold, and let $B$ be a globally generated holomorphic vector bundle of rank $n$ on $X$. As in the case of line bundles, $G$ induces a map $\Phi_B$ to $\gr(V,n)$, where $V=H^0(X,B)^*$, as follows

\begin{align*}
    \Phi_B: X&\to \gr(V,n)\\
    p&\mapsto (\ker v_p)^\circ
\end{align*}
where $v_p:H^0(X,B)\to B_p$ is the evaluation map at the point $p\in X$, and $(\ker v_p)^\circ$ denotes the annihilator of $\ker v_p$. Note that the map is well defined since $B$ is globally generated. As a last general remark, consider the following composition 
$$X\xrightarrow{\Phi_B}\gr(V,n)\xhookrightarrow{P} \bP(\Lambda^n V)$$ where the last map is the Pl\"ucker embedding, i.e. the embedding induced from the line bundle $\det Q$. It holds that $\Phi_B^*Q\cong B$, hence $\det B\cong(\Phi_B\circ P)^*\cO_{\bP(\Lambda^nV)}(1)$. 
From this observation, we obtain the following:

\begin{lemma}\label{lemma_vector_bundle_to_determinat_line_bundle}
Let $X$ be a compact complex manifold, and let $B$ be a globally generated vector bundle on X, then $\det B$ is base point free. Furthermore, if $\Phi_B$ is an embedding, then $\det B$ is very ample.
\end{lemma}

Returning to the Hilbert schemes, let $S$ be a complex compact surface and let $L$ be a $(n-1)$-very ample line bundle on $S$. Proposition \ref{prop_tautological_globally_generated} guarantees that the vector bundle $\hilb{L}$ is globally generated, hence we have the well defined map 
\begin{align*}
    \Phi_{\hilb{L}}:\hilb{S}&\to \gr(H^0(\hilb{S},\hilb{L})^*,n).\\
    \eta &\mapsto (\ker v_\eta)^\circ
\end{align*}


In this case, we have another map into a Grassmannian. As a matter of fact, recall that $(n-1)-$generation of $L$ implies that for all $\eta\in \hilb{S}$ the restriction $r_\eta :H^0(S,L)\to H^0(\eta,L\otimes\cO_\eta)$ is surjective, in particular this allows us to define a map 

\begin{align*}
    \varphi_{\hilb{L}}:\hilb{S}&\to \gr(H^0(S,L)^*,n).\\
    \eta &\mapsto (\ker r_\eta)^\circ 
\end{align*}

In \cite{Catanese_Goettsche}, the following result is proven:

\begin{theorem}
The morphism $\varphi_{\hilb{L}}:\hilb{S}\to \gr(H^0(S,L)^*,n)$ is an embedding if and only if $L$ is $n$-very ample.
\end{theorem}



Consider now the map 
\begin{align*}
    g:H^0(S,L)\to H^0(\hilb{S},\hilb{L})
\end{align*} 
where $g(s)(\eta)=s|_\eta\in H^0(\eta,\cO_\eta\otimes L)$ for all $\eta\in \hilb{S}$. The map $g$ is injective and, thanks to Proposition \ref{prop_zeroth_cohom_F^[n]}, it is an isomorphism, hence it induces an isomorphism $$\tilde{g} :\gr(H^0(S,L)^*,n)\to \gr(H^0(\hilb{S},\hilb{L})^*,n).$$ Observe now that $r_\eta=v_\eta\circ g$ for all $\eta\in\hilb{S}$, hence the following diagram is commutative 

% https://q.uiver.app/#q=WzAsNCxbMSwwLCJcXGdyKEheMChTLEwpXiosbikiXSxbMSwyXSxbMSwxLCJcXGdyKEheMChcXGhpbGJ7U30sXFxoaWxie0x9KV4qLG4pIl0sWzAsMCwiXFxoaWxie1N9Il0sWzAsMiwiXFx0aWxkZXtnfSJdLFszLDAsIlxcdmFycGhpX3tMfSJdLFszLDIsIlxcUGhpX3tcXGhpbGJ7TH19IiwyXV0=
\[\begin{tikzcd}
	{\hilb{S}} & {\gr(H^0(S,L)^*,n)} \\
	& {\gr(H^0(\hilb{S},\hilb{L})^*,n)} \\
	& {}
	\arrow["{\varphi_{\hilb{L}}}", from=1-1, to=1-2]
	\arrow["{\Phi_{\hilb{L}}}"', from=1-1, to=2-2]
	\arrow["{\tilde{g}}", from=1-2, to=2-2]
\end{tikzcd}\]

that is $\Phi_{\hilb{L}}=\tilde{g}\circ\varphi_{\hilb{L}}$.


\begin{proposition}\label{prop_det_of_tautological_bpf_va}
    Let $S$ be a compact complex smooth surface and let $L$ line bundle on $S$, the following hold:
    \begin{itemize}
        \item If $L$ is $(n-1)$-very ample then $\det \hilb{L}$ is base point free;
        \item If $L$ is $n$-very ample then $\det \hilb{L}$ is very ample.
    \end{itemize}
    In both cases the same property is shared with the line bundle $L_n-\delta$.
\end{proposition}

\begin{proof}
If $L$ is $(n-1)$-very ample, Proposition \ref{prop_tautological_globally_generated} implies that $\hilb{L}$ is globally generated, then Lemma \ref{lemma_vector_bundle_to_determinat_line_bundle} implies that $\det \hilb{L}$ is base point free as wanted.

Suppose now that $L$ is $n$-very ample. Under this assumption, the map $\Phi_{\hilb{L}}=\tilde{g}\circ\varphi_{\hilb{L}}$ is a composition of two embeddings, hence an embedding itself. We can conclude using Lemma \ref{lemma_vector_bundle_to_determinat_line_bundle}.

The last observation follows immediately from Proposition \ref{prop_det_of_tautological}.
\end{proof}



\section{Hyperk\"ahler manifolds}\label{sect_hyper}

Let $X$ be a K\"ahler manifold, we say that $X$ is \emph{hyperk\"ahler} (or \emph{irreducible holomorphic symplectic}) if it is simply connected and if $H^0(X,\Omega^2_X)$ is generated by an everywhere non-degenerate holomorphic 2-form.
%%%%%%%%%%%%%%%%%%%%%%%%%%%%%%%%%%%%%%%%%%%%%%%%%%%%%%%
\begin{comment}
   \marginpar{O'Grady's examples can also be deformed in families, but they occur only in dimension 6 and 10, maybe mention this} 
\end{comment}
%%%%%%%%%%%%%%%%%%%%%%%%%%%%%%%%%%%%%%%%%%%%%%%%%%%%%%%
There are two known families of examples that occur in every possible dimension:
\begin{itemize}
    \item Given a $K3$ surface $S$ its Hilbert scheme on $n$ points $\hilb[n]{S}$ is a hyperk\"ahler manifold;
    \item  Given an abelian surface $T$ consider the Hilbert scheme $\hilb[n+1]{T}$. There is a natural map $\hilb[n+1]{T}\to T$ given by the sum. The fiber over $0$ is a hyperk\"ahler manifold called \emph{generalized Kummer variety}, denoted by $\Kum^n(T)$.
\end{itemize}

Furthermore, deformations of hyperk\"ahler manifolds are still hyperk\"ahler. In particular, deformations of $\hilb[n]{S}$ are called of $\hilb[n]{K3}$-type, while deformations of $\Kum^n(T)$ are called of $\Kum^n$-type.
The only other known examples of hyperk\"ahler manifolds are O'Grady's "exceptional" examples, namely OG6 and OG10, which occur only in dimensions 6 and 10, respectively. 

Given an hyperk\"ahler manifold, the second cohomology group $H^2(X,\bZ)$ can be endowed with a quadratic form $q$ called the Beauville-Bogomolov-Fujiki form (or BFF form), so that $(H^2(X,\bZ),q)$ has the structure of a lattice (see \cite{huybrechts1997compacthyperkaehlermanifoldsbasic} for a reference). In this regard, let us recall some definitions. Let $(\Lambda,q)$ a lattice:
\begin{itemize}
    \item we say an element $\alpha$ is \emph{primitive} if $\alpha=k\alpha'$ implies $k=\pm1$;
    \item The \emph{multiplicity} of an element $\alpha$ is $m$ if there exists a primitive element $\alpha'$ such that $\alpha=m\alpha'$;
    \item The \emph{divisibility} of an element $\alpha$, denoted by $\Div(\alpha)$, is the multiplicity of $(\cdot,\alpha)_q$ as an element of the dual lattice $\Lambda^\vee$. Equivalently $\Div(\alpha)=m$ if $m$ generates the ideal $\{(\alpha',\alpha)_q|\alpha'\in\Lambda\}$;
    \item Let $\{e_i\}$ be a $\bZ$-basis of the lattice $\Lambda$. The \emph{discriminant} of the $\Lambda$ is defined as the determinant of the matrix $(q(e_i,e_j))_{i,j}$.
\end{itemize}

We report a characterization of projective hyperk\"ahler manifolds.
\begin{theorem}[\cite{huybrechts1997compacthyperkaehlermanifoldsbasic}, Theorem 3.11]\label{theo_projective_hyperkahler_manifolds}
    Let $X$ be a hyperk\"ahler manifold. Then $X$ is projective if and only if there exists a line bundle $L$ on $X$ with $q(L)>0$.
\end{theorem}

The lattice structure of the $\hilb[n]{K3}$-type and $\Kum^n$-type varieties is described in the following proposition. Fix the notation $\Lambda_{K3}\cong U^{\oplus 3}\oplus E_8(-1)^{\oplus 2}$ for the $K3$ lattice, where $U$ is the standard hyperbolic lattice and $E_8(-1)$ is the unimodular root-lattice $E_8$ change by a sign.

\begin{proposition}\label{p_k3_type_lattice}
    Fix a $K3$ surface $S$. Then the following holds:
    \begin{itemize}
        \item[(i)] For every $n\geq 2$, there is an orthogonal decomposition of lattices $$H^2(\hilb[n]{S},\bZ)\cong H^2(S,\bZ)\oplus \bZ\delta,$$
        where $\delta$ is an integral $(1,1)$-class with $q(\delta)=-2(n-1)$. The class $2\delta$ is represented by the Hilbert-Chow divisor.
        \item[(ii)] There is an orthogonal decomposition:
        $$\Pic(\hilb[n]{S})\cong\Pic(S)\oplus\bZ\delta.$$
        \item[(iii)]For each irreducible symplectic variety $X$ of $\hilb[n]{K3}$-type there exists an isomorphism $$H^2(X,\bZ)\cong\Lambda_{K3}\oplus \bZ\delta.$$  
    \end{itemize}
\end{proposition}

\begin{proposition}\label{p_kum_type_lattice}
    Fix an abelian surface $T$. Then the following holds:
    \begin{itemize}
        \item[(i)] For every $n\geq 2$, there is an orthogonal decomposition of lattices $$H^2(\Kum^n(T),\bZ)\cong H^2(T,\bZ)\oplus \bZ\delta,$$
        where $\delta$ is an integral $(1,1)$-class with $q(\delta)=-2(n+1)$. The class $2\delta$ is represented by the Hilbert-Chow divisor.
        \item[(ii)] There is an orthogonal decomposition:
        $$\Pic(\Kum^n(T))\cong\Pic(T)\oplus\bZ\delta.$$
        \item[(iii)]For each irreducible symplectic variety $X$ of $\Kum^n$-type there exists an isomorphism $$H^2(X,\bZ)\cong H^2(T,\bZ)\oplus \bZ\delta$$
        for some abelian variety $T$.
    \end{itemize}
\end{proposition}

\begin{remark}\label{remark_how_to_compute_divisibility}
Proposition \ref{p_k3_type_lattice} implies that given any element $\alpha\in H^2(X,\bZ)$ where $X$ is a hyperk\"ahler manifold of $\hilb[n]{K3}$-type, then $\alpha=a\lambda+b\delta$ with $\lambda\in \Lambda_{K3}$ a primitive element. In particular, recalling that $\Lambda_{K3}$ is a unimodular lattice we get $$\Div(\alpha)=\gcd(a,2b(n-1)).$$ 
Analogously for $X$ $\Kum^n$-type we have $$\Div(\alpha)=\gcd(a,2b(n+1)),$$
since $H^2(T,\bZ)$ is unimodular for any abelian surface $T$.
\end{remark}





\section{Moduli space of polarized hyperk\"ahler manifolds}\label{sect_moduli}
We denote by $\Sigma_{d,t}^n$ the moduli space parameterizing pairs $(X,L)$ with $X$ a hyperk\"ahler manifold of $\hilb{K3}$-type and $L$ a primitive ample line bundle with $q(L)=2d$ and $\Div(L)=t$. We used the symbol $\Upsilon_{d,t}^n$ to denote the analogous moduli space for $\Kum^n$-type manifolds.

In \cite{zbMATH06387271} and \cite{onorati2022connectedcomponentsmodulispaces}, the authors give a characterization of the number of connected components of $\Sigma_{d,t}^n$ and $\Upsilon_{d,t}^n$, respectively. In the following, write $\phi$ for the Euler function, and let $\rho(a)$ denote the number of distinct primes in the factorization of a positive integer $a$.



\begin{theorem}[\cite{zbMATH06387271}, Prop 3.1]\label{theo_count_connected_components_hilb_k3}
    Let $t$ be a divisor of $(2d, 2n-2)$, and set 
    \begin{equation}
\begin{array}{lll}
\tilde{d}=\frac{2d}{\gcd(2d,2n-2)}, & \tilde{n}=\frac{2n-2}{\gcd(2d,2n-2)}, & g=\frac{\gcd(2d,2n-2)}{t} \\
 & & \\
w=\gcd(g,t), & g_1=\frac{g}{w}, & t_1=\frac{t}{w}.
\end{array}
\end{equation}
Put $w=w_+(t_1)w_-(t_1)$ where $w_+(t_1)$ is the product of all powers of primes dividing $\gcd(w, t_1)$.

     \begin{itemize}
       \item $|\Sigma_{d,t}^n|=w_+(t_1)\phi(w_-(t_1))2^{\rho(t_1)-1}$ if $t>2$ and one of the following sets of conditions hold: 
       \begin{itemize}
           \item[(i)] $g_1$ is even, $(\tilde{d}, t_1)=(\tilde{n}, t_1)=1$ and the residue class $-\tilde{d}/\tilde{n}$ is a quadratic residue modulo $t_1$;
           \item[(ii)] $g_1, t_1$, and $\tilde{d}$ are odd, $(\tilde{d}, t_1)=(\tilde{n}, 2t_1)=1$ and $-\tilde{d}/\tilde{n}$ is a quadratic residue modulo $2t_1$;
           \item[(iii)] $g_1, t_1$, and $w$ are odd, $\tilde{d}$ is even, $(\tilde{d}, t_1)=(\tilde{n}, 2t_1)=1$ and $-\tilde{d}/(4\tilde{n})$ is a quadratic residue modulo $t_1$;
       \end{itemize}
       \item $|\Sigma_{d,t}n|=w_+(t_1)\phi(w_-(t_1))2^{\rho(t_1/2)-1}$ if $t>2$, $g_1$ is odd, $t_1$ is even, $(\tilde{d}, t_1)=(\tilde{n}, 2t_1)=1$ and $-\tilde{d}/\tilde{n}$ is a quadratic residue modulo $2t_1$;
       \item $|\Sigma_{d,t}^n|=1$ if $t\leq 2$ and one of the following sets of conditions hold: 
       \begin{itemize}
           \item[(i)] $g_1$ is even, $(\tilde{d}, t_1)=(\tilde{n}, t_1)=1$ and the residue class $-\tilde{d}/\tilde{n}$ is a quadratic residue modulo $t_1$;
           \item[(ii)] $g_1, t_1$, and $\tilde{d}$ are odd, $(\tilde{d}, t_1)=(\tilde{n}, 2t_1)=1$ and $-\tilde{d}/\tilde{n}$ is a quadratic residue modulo $2t_1$; 
           \item[(iii)] $g_1, t_1$, and $w$ are odd, $\tilde{d}$ is even, $(\tilde{d}, t_1)=(\tilde{n}, 2t_1)=1$ and $-\tilde{d}/(4\tilde{n})$ is a quadratic residue modulo $t_1$;
           \item[(iv)] $g_1$ is odd, $t_1$ is even, $(\tilde{d}, t_1)=(\tilde{n}, 2t_1)=1$ and $-\tilde{d}/\tilde{n}$ is a quadratic residue modulo $2t_1$;
       \end{itemize}
       \item $|\Sigma_{d,t}^n|=0$, else.
     \end{itemize}
\end{theorem}



\begin{theorem}[\cite{onorati2022connectedcomponentsmodulispaces}, Thm 5.3]\label{theo_count_connected_components_kumm_type}
Let $t$ be a divisor of $\gcd(2d,2n+2)$. Consider the following notations,
\begin{equation}
\begin{array}{lll}
d_1=\frac{2d}{\gcd(2d,2n+2)}, & n_1=\frac{2n+2}{\gcd(2d,2n+2)}, & g=\frac{\gcd(2d,2n+2)}{t} \\
 & & \\
w=\gcd(g,t), & g_1=\frac{g}{w}, & t_1=\frac{t}{w}.
\end{array}
\end{equation}
For $w$ and $t_1$ as defined above, we write $w=w_+(t_1)w_-(t_1)$, where $w_+(t_1)$ is the product of all powers of the primes that appear in the factorization of $w$ and that divide $\gcd(w,t_1)$. We have:
\begin{itemize}

\item $|\Upsilon_{d,t}^n|=w_+(t_1)\phi(w_-(t_1))2^{\rho(t_1)-1}$ if $t>2$ and one of the following holds:
 \begin{itemize}
 \item $g_1$ is even, $\gcd(d_1,t_1)=1=\gcd(n_1,t_1)$ and $-d_1/n_1$ is a quadratic residue mod~$t_1$;
 \item $g_1,t_1$ and $d_1$ are odd, $\gcd(d_1,t_1)=1=\gcd(n_1,2t_1)$ and $-d_1/n_1$ is a quadratic residue mod~$2t_1$;
 \item $g_1,t_1$ and $w$ are odd, $d_1$ is even, $\gcd(d_1,t_1)=1=\gcd(n_1,2t_1)$ and $-d_1/4n_1$ is a quadratic residue mod~$t_1$.
 \end{itemize}

\item $|\Upsilon_{d,t}^n|=w_+(t_1)\phi(w_-(t_1))2^{\rho(t_1/2)-1}$ if $t>2$, $g_1$ is odd, $t_1$ is even, $\gcd(d_1,t_1)=1=\gcd(n_1,2t_1)$ and $-d_1/n_1$ is a quadratic residue mod~$2t_1$.

\item $|\Upsilon_{d,t}^n|=1$ if $t\leq2$ and one of the following holds:
\begin{itemize}
 \item $g_1$ is even, $\gcd(d_1,t_1)=1=\gcd(n_1,t_1)$ and $-d_1/n_1$ is a quadratic residue mod~$t_1$;
 \item $g_1,t_1$ and $d_1$ are odd, $\gcd(d_1,t_1)=1=\gcd(n_1,2t_1)$ and $-d_1/n_1$ is a quadratic residue mod~$2t_1$;
 \item $g_1,t_1$ and $w$ are odd, $d_1$ is even, $\gcd(d_1,t_1)=1=\gcd(n_1,2t_1)$ and $-d_1/4n_1$ is a quadratic residue mod~$t_1$;
 \item $g_1$ is odd, $d_1$ is even, $\gcd(d_1,t_1)=1=\gcd(n_1,2t_1)$ and $-d_1/n_1$ is a quadratic residue mod~$2t_1$.
 \end{itemize}

\item $|\Upsilon_{d,t}^n|=0$ otherwise.
\end{itemize}
\end{theorem}

\begin{proposition}[\cite{huybrechts1997compacthyperkaehlermanifoldsbasic}, 1.14]\label{prop_irreducibility_connected_components}
    For any $n,d,t$ positive integers with $n>1$ the moduli space $\Sigma_{d,t}^n$ and $\Upsilon_{d,t}^n$ is smooth. In particular every connected component of these moduli spaces is irreducible.
\end{proposition}


Being interested in studying the base point freeness and very ampleness of general elements of these moduli spaces, we now investigate the openness condition of these properties. Following the notation in \cite{nugent2025moduliamplelinebundles}, let $f:X\to Y$ be a morphism and $L$ be a line bundle on $X$, define $$\cB\cF(f,L):\left\{y\in Y: L_y \text{ is base point free}\right\}$$
and 
$$\cV\cA(f,L):\left\{y\in Y: L_y \text{ is very ample}\right\}.$$

Then the following theorem holds.

\begin{theorem}[\cite{nugent2025moduliamplelinebundles}, Lemma 7.3 - Thm 7.4]\label{theorem_openness_bpf_va}
    Let $f:X\to Y$ be a proper flat morphism of locally noetherian scheme and let L be a line bundle on $X$ such that $h^0(X_y,L_y)$ is constant for all $y\in Y$. Then $\cB\cF(f,L)$ and $\cV\cA(f,L)$ are open. 
\end{theorem}

Specilizing to hyperk\"ahler manifolds we obtain the following.

\begin{corollary}\label{cor_openess_of_bpf_va}
Let $f:X\to S$ be a family of hyperk\"ahler manifolds over a connected base $S$. Let $L$ be a line bundle on $X$ such that $q(L_s)>0$ for all $s\in S$. Then:
\begin{itemize}
    \item if $L_0$ is base point free for some $0\in S$ then $L_s$ is base point for all $s$ in an open neighborhood of $0\in S$;
    \item if $L_0$ is very ample for some $0\in S$ then $L_s$ is very ample for all $s$ in an open neighborhood of $0\in S$;
\end{itemize}
\end{corollary}
\begin{proof}
    We follow the proof of Proposition 2.6 in \cite{Varesco_2023} which we report here for convenience. We consider the situation where $L_0$ is very ample as the other follows analogously. Since $q(L_0)>0$ then $L_0$ is big and nef by the Beauville-Fujiki relations [\cite{Gross2003CalabiYau}, Prop. 23.14], hence $h^i(X_0,L_0)=0$ for $i>0$ thanks to the Kodaira vanishing theorem. The semicontinuity theorem then implies that $h^i(X_s,L_s)=0$ for all $i>0$ for all $s$ in a neighborhood $U$ of $0\in S$. By flatness of $f$, the Euler characteristic $\chi(L_s)$ does not depend on $s$, therefore $h^0(X_s,L_s)$ is constant for all $s\in U$. By theorem \ref{theorem_openness_bpf_va}, the set $\cV\cA(f,L)$ is open in $U$ hence in $S$ as wanted.
\end{proof}


\begin{lemma}\label{lemma_bfp_does_not_require_ampleness}
Let $d,n,t$ be positive integers with $n\geq 2$, let $X$ be an hyperk\"ahler manifold and let $L\in \Pic(X)$ be base point free such that $q(L)=2d$, $\Div(L)=t$. Then the general element in a connected component of $\Sigma_{d,n}^t$, where $\dim_{\bC}(X)=n$,  has base point free polarization. The same holds for $\Upsilon_{d,n}^t$.
\end{lemma}
\begin{proof}
    A very general deformation of $(X,L)$, call it $(X_A,L_A)$, has Picard rank equal to 1 hence, by the Projectivity criteria in Theorem \ref{theo_projective_hyperkahler_manifolds}, the line bundle $L_A$ is ample since $q(L_A)>0$. The properties of being ample and base point free are open conditions hence, we can find a generic deformation $(X_{AF},L_{AF})$ such that $L_{AF}$ is ample, base point free and satisfy the same numerical properties of $L$. Consequently, the element $(X_{AF},L_{AF})$ defines an element of $\Sigma_{d,n}^t$. We can conclude that the generic element of $\Sigma_{d,n}^t$, in the same connected component as $(X,L)$, has base point free polarization since the property is open in families by Corollary \ref{cor_openess_of_bpf_va} and the connected components are irreducible by Proposition \ref{prop_irreducibility_connected_components}.
\end{proof}


\section{Base point freeness and very ampleness of a general element}\label{sect_main_theorems}

For K3 and Abelian surfaces, there is a characterization for $k$-very ampleness that can be applied in the case the Picard rank is 1.

\begin{proposition}[\cite{debarre2020hyperkahlermanifolds}, Cor. 3.8]\label{prop_k_very_ample_k3}
Let $S$ be a K3 surface such that $\Pic(S)\cong\bZ L$ with $L$ ample and $L^2=2e$. Then the line bundle $L^{\otimes a}$ is $k$-very ample if and only if either
\begin{itemize}
    \item $a=1$ and $k\leq \frac{e}{2}$;
    \item $a\geq 2$ and $k\leq 2(a-1)e-2$.
\end{itemize}
\end{proposition}

\begin{proposition}[\cite{Alagal_2016}, Theo. 4.3 - Prop. 1.5]\label{prop_k_very_ample_abelian}
Let $T$ be an Abelian surface such that $\Pic(T)\cong\bZ L$ with $L$ ample and $L^2=2e$. Then the line bundle $L^{\otimes a}$ is $k$-very ample if and only if either
\begin{itemize}
    \item $a=1$ and $k\leq \lfloor\frac{e-3}{2}\rfloor$;
    \item $a\geq 2$ and $k\leq 2(a-1)e-2$.
\end{itemize}
\end{proposition}

We can now apply Proposition \ref{prop_det_of_tautological_bpf_va} to obtain the following corollaries.

\begin{corollary}\label{cor_bpf_va_line_bundle_for_k3}
Let $S$ be a K3 surface such that $\Pic(S)\cong\bZ L$ with $L$ ample and $L^2=2e$. Then the line bundle $aL_n-\delta$ on $\hilb{S}$ is base point free if either
\begin{itemize}
    \item $a=1$ and $n\leq \frac{e}{2}+1$;
    \item $a\geq 2$ and $n\leq 2(a-1)e-1$.
\end{itemize}
Also, the line bundle $aL_n-\delta$ on $\hilb{S}$ is very ample if either
\begin{itemize}
    \item $a=1$ and $n\leq \frac{e}{2}$;
    \item $a\geq 2$ and $n\leq 2(a-1)e-2$.
\end{itemize}
\end{corollary}

\begin{corollary}\label{cor_bpf_va_line_bundle_for_abelian}
Let $T$ be an Abelian surface such that $\Pic(T)\cong\bZ L$ with $L$ ample and $L^2=2e$. Then the line bundle $aL_{n+1}-\delta$ on $\Kum^n(T)$ is base point free if either
\begin{itemize}
    \item $a=1$ and $n\leq \lfloor\frac{e-3}{2}\rfloor$;
    \item $a\geq 2$ and $n\leq 2(a-1)e-2$.
\end{itemize}
Also, the line bundle $aL_{n+1}-\delta$ on $\Kum^n(T)$ is very ample if either
\begin{itemize}
    \item $a=1$ and $n\leq \lfloor\frac{e-3}{2}\rfloor-1$;
    \item $a\geq 2$ and $n\leq 2(a-1)e-3$.
\end{itemize}
\end{corollary}

\begin{remark}
    Under a little abuse of notation, we are using $L_{n+1}$ to denote the line bundle on $\hilb[n+1]{T}$ and its restriction to $\Kum^n(T)$. This is not a problem since both very ampleness and base point freeness behave nicely under restriction.
\end{remark}

\begin{theorem}
    
    Let $n,d$ be positive integers with $n\geq 2$.\\
    Let $(X,H)$ be a general element of $\Sigma_{d,n}^1$ then 
    \begin{itemize}
        \item the line bundle $H$ is base point free;
        \item if $d\geq n+1$ the line bundle $H$ is very ample.
    \end{itemize}
    Let $(X,H)$ be a general element of $\Upsilon_{d,n}^1$ then
    \begin{itemize}
        \item if $d\geq 3$ the line bundle $H$ is base point free;
        \item if $d\geq n+4$ the line bundle $H$ is very ample.
    \end{itemize}
\end{theorem}

\begin{proof}
    Let us present the proof for the case of $\Sigma_{d,n}^1$ as the other case follows analogously. Let $S$ be a K3 surface such that $\Pic(S)\cong \bZ L$ with $L$ ample and $L^2=2d$. Recalling that $0$-very ample is the same as base point free, Corollary \ref{cor_bpf_va_line_bundle_for_k3} implies that $L$ is base point free. Then according to Lemma \ref{lemma_bpf_of_L_n_on_Hilb_n}, the line bundle $L_n\in \Pic(\hilb{S})$ is also base point free. Using Proposition \ref{p_k3_type_lattice} it is primitive and satisfies $q(L_n)=2d$, $\Div(L_n)=1$. Observing that $\Sigma_{d,n}^1$ is connected by Theorem \ref{theo_count_connected_components_hilb_k3}, we can conclude thanks to Lemma \ref{lemma_bfp_does_not_require_ampleness}.\\
    The proof of the very ampleness part follows the same proof as in Theorem \ref{theo_general_in_connected_component} combined once again with the irreducibility of $\Sigma_{d,n}^1$.
\end{proof}

\begin{remark}
    The discrepancy between the $\hilb{K3}$-type and the $\Kum^n$-type, for the base point freeness, is due to the existence of Abelian surfaces with ample line bundle $L$ of low self intersection, namely $L^2=2$ and $L^2=4$, which are not base point free. However, using the automorphisms induced by the group structure of Abelian surfaces, it can be proven that, even though $L$ on $T$ is not base point free, the line bundle $L_{n+1}$ on $\Kum^n(T)$ is. This allows us to obtain a result like the $\hilb{K3}$-type case, see \cite{Varesco_2023} Theorem 3.1.
\end{remark}

\begin{theorem}\label{theo_general_in_connected_component}
    Let $n,d,t$ be positive integers with $t,n\geq 2$, let $\tau =\frac{t^2}{2(t-1)}$ then
    \begin{itemize}
        \item if $d\geq (\tau-1)n+\tau+1$ there exists a connected component of $\Sigma_{d,n}^t$ such that the polarization of a general element $(X,H)$ is base point free;
        \item if $d\geq (\tau-1)n+2\tau+1$ there exists a connected component of $\Sigma_{d,n}^t$ such that the polarization of a general element $(X,H)$ is very ample;
        \item if $d\geq (\tau-1)n+2\tau-1$ there exists a connected component of $\Upsilon_{d,n}^t$ such that the polarization of a general element $(X,H)$ is base point free;
        \item if $d\geq (\tau-1)n+3\tau-1$ there exists a connected component of $\Upsilon_{d,n}^t$ such that the polarization of a general element $(X,H)$ is very ample;
    \end{itemize}
These hold unless the considered moduli space is empty.
\end{theorem}

\begin{proof}
    Let us focus on $\Upsilon_{d,n}^t$ in the case of very ampleness as the proofs of the other instances are very similar. Fix $t,n$ positive integers with $t,n\geq 2$ such that $t|2(n+1)$ (this is a necessary condition for $t$ to be a possible divisibility) and pick a polarized Abelian surface $(T,L)$ with $\Pic(T)\cong \bZ L$ and $L^2=e$. According to Remark \ref{remark_how_to_compute_divisibility} it then holds $$\Div(tL_{n+1}-\delta)=(t,2(n+1))=t$$
    and $$q(tL_{n+1}-\delta)=2\left(t^2e-n-1\right)=:2d.$$
By Corollary \ref{cor_bpf_va_line_bundle_for_abelian}, the line bundle $tL_{n+1}-\delta$ is very ample if $$d\geq (\tau-1)n+3\tau-1.$$
Thanks to Corollary \ref{cor_openess_of_bpf_va} and Proposition \ref{prop_irreducibility_connected_components}, the polarization of a general element $(X,H)$, in the connected component of $\Upsilon_{d,n}^t$ containing $(\Kum^n(T),tL_{n+1}-\delta)$, is very ample.
    \end{proof}

\begin{corollary}
        Let $n,d$ be positive integers with $n\geq 2$.\\
        Let $(X,H)$ be a general element of $\Sigma_{d,n}^t$ and suppose $t=2$ or $t=p^a$ with $p> 2$ a prime number such that $p^{a+1}$ does not divide $(2d,2(n-1))$ then
    \begin{itemize}
        \item if $d\geq (\tau-1)n+\tau+1$ the line bundle $H$ is base point free;
        \item if $d\geq (\tau-1)n+2\tau+1$ the line bundle $H$ is very ample;
    \end{itemize}
    Let $(X,H)$ be a general element of $\Upsilon_{d,n}^t$ and suppose $t=2$ or or $t=p^a$ with $p> 2$ a prime number such that $p^{a+1}$ does not divide $(2d,2(n+1))$ then
    \begin{itemize}
        \item if $d\geq (\tau-1)n+2\tau-1$ the line bundle $H$ is base point free;
        \item if $d\geq (\tau-1)n+3\tau-1$ the line bundle $H$ is very ample;
    \end{itemize}
These hold unless the considered moduli space is empty.
\end{corollary}

\begin{proof}
    This follows from Theorem \ref{theo_count_connected_components_hilb_k3} or Theorem \ref{theo_count_connected_components_kumm_type} after checking that these moduli spaces have one connected component. For the sake of it, we are going to check that this works for $\Sigma_{d,n}^t$ with $t=p^a$ power of a prime number greater than $2$ such that $p^{a+1}$ does not divide $(2d,2(n-1))$. Under these conditions, using the same notation as in Theorem \ref{theo_count_connected_components_hilb_k3}, we have that the prime $p$ does not divide $g$ hence, since $t=p^a$, it holds $w=1$ and $t_1=t$. Therefore, we have that $t_1$ is odd, hence we necessarily are in the first case of Theorem \ref{theo_count_connected_components_hilb_k3}. Thus $\rho(t_1)=\rho(t)=1$ and $w_+(t_1)=w_-(t_1)=1$ hence $|\Sigma_{d,n}^t|=1$, as desired.
\end{proof}



Note that a lot of the moduli spaces that do not fall in the conditions of Theorem \ref{theo_general_in_connected_component} are actually empty. For example, from \ref{remark_how_to_compute_divisibility} it follows that a necessary condition for $\Sigma_{d,n}^t$ to not be empty is that $t|2(n-1)$. With this in mind, we are giving a necessary and sufficient condition for $\Sigma_{d,n}^t$ and $\Upsilon_{d,n}^t$ to be non empty.

\begin{proposition}\label{prop_char_of_non_emptyness}
    Let $n,d,t$ be integers with $n\geq 2$ then
    \begin{equation*}\Sigma_{d,n}^t\neq \emptyset \iff 
        \begin{cases}
            t|(2n-2)\\
            d\equiv_{t^2}-b^2(n-1) \text{ for some } b \text{ such that } $(b,t)=1$
        \end{cases}
    \end{equation*}
    and 
    \begin{equation*}\Upsilon_{d,n}^t\neq \emptyset \iff 
        \begin{cases}
            t|(2n+2)\\
            d\equiv_{t^2}-b^2(n+1) \text{ for some } b \text{ such that } $(b,t)=1$
        \end{cases}
    \end{equation*}
\end{proposition}
\begin{proof}
    Suppose $\Sigma_{d,n}^t$ is non empty, then we can pick an element $(X,H)$. Thanks to Proposition \ref{p_k3_type_lattice}, we can find a K3 surface $S$ and a line bundle $L$ on $S$ such that $H\cong aL+b\delta$ for some integers $a,b$. Recalling that $q(\delta)=-2(n-1)$, the line bundle $aL+b\delta$ satisfies the following numerical properties
    \begin{align*}
        a^2q(L)-2b^2(n-1)&=2d,\\
        (a,2b(n-1))&=t,\\
        (a,b)&=1.
    \end{align*}
    From the second equation, we immediately obtain $t|(2n-2)$ as wanted. From the first one we can write $$d\equiv_{t^2}-b^2(n-1),$$ and since $t|a$ and $(a,b)=1$ it holds $(t,b)=1$, as wanted.\\
    Suppose now that we have positive integers $d,n,t$ with $n\geq2 $ satisfying 
    \begin{equation*}
       \begin{cases}
            t|(2n-2)\\
            d\equiv_{t^2}-b^2(n-1) \text{ for some } b \text{ such that } $(b,t)=1$.
        \end{cases}
    \end{equation*}
    Observe that the number $$e:=\frac{d+b^2(n-1)}{t^2}$$ is a positive integer, so there exists a K3 surface $S$ with line bundle $L$ such that $L^2=2e$, and after deforming we can assume $\Pic(S)\equiv \bZ L$, so $L$ is primitive. Consequently,  $tL_n+b\delta$ is primitive as $(t,b)=1$, and it satisfies 
    \begin{align*}
        q(tL_n+b\delta)&=2t^2e-2b^2(n-1)=2d,\\
        \Div(tL_n+b\delta)&=(t,2b(n-1))=t.
    \end{align*}
    For general deformation $(X,L)$ of $(\hilb{S},tL_n-b\delta)$, the Picard group is generated by $H$, hence by the projectivity criteria in Theorem \ref{theo_projective_hyperkahler_manifolds}, the line bundle $H$ or $-H$ is ample. Suppose $H$ is, then $(X,H)\in \Sigma_{d,n}^t$ as wanted. If $-H$ is ample, we can conclude analogously, since $\Div(-H)=\Div(H)$ and $q(H)=q(-H)$.
    The proof for $\Upsilon_{d,n}^t$ similarly follows.
\end{proof}

We would like to conclude with one more remark.\\
Recall the following notion from \cite{Lazarsfeld2004PositivityI} (Definition 1.8.4).

Let $X$ be a projective variety and $B$ an ample line bundle on $X$ that is generated by its global sections. A coherent sheaf $\mathcal{F}$ on $X$ is \emph{$m$-regular with respect to $B$} if $$H^i(X,\mathcal{F}\otimes B^{\otimes(m-i)})=0$$ for $i>0$.

What we are interested in is the following proposition:
\begin{proposition}[\cite{Lazarsfeld2004PositivityI}, Prop 1.8.22]\label{prop_0_regular_gives_very_ample}
Let $X$ be an irreducible projective variety of dimension $n$ and $B$ an ample line bundle that is generated by its global sections and let $N$ be a line bundle that is 0-regular with respect to $B$. Then $N\otimes B$ is very ample.
\end{proposition} 

From which we obtain the following lemma.

\begin{lemma}
    Let $X$ be a hyperk\"ahler manifold of complex dimension $n$ and let $L$ be an ample base point free line bundle on $X$. Then $L^{n+2}$ is very ample.
\end{lemma}

\begin{proof}
    If we can apply Proposition \ref{prop_0_regular_gives_very_ample} with $B=L$ and $N=L^{n+1}$, then $N\otimes B=L^{n+2}$ is very ample and we are done. We need to show that $$H^i(X,N\otimes B^{\otimes(0-i)})=0\quad \forall i>0$$
    which, with our choices of $N$ and $B$, becomes $$H^i(X,L^{\otimes(n+1-i)})=0\quad \forall i>0.$$
    If $i<n+1$, then Kodaira vanishing can be applied.
    If $i\geq n+1$, then the cohomology group is trivial since $h^{k,l}(X,H)=0$ for any line bundle $H\in \Pic(X)$, if $k>\dim_{\mathbb{C}} X=n$.
\end{proof}

\begin{corollary}
     Let $n,d,t$ be positive integers with $n\geq 2$, let $\tau =\frac{t^2}{2(t-1)}$. Then if $(X,H)$ is a general element of $\Sigma_{d,n}^1$ or $\Upsilon_{d,n}^1$ then $H$ is base point free and $H^{n+2}$ is very ample. If $t\geq 2$ the following also holds:
    \begin{itemize}
        \item if $d\geq (\tau-1)n+\tau+1$ there exists a connected component of $\Sigma_{d,n}^t$ such that the polarization of a general element $(X,H)$ is base point free and $H^{n+2}$ is very ample;
        \item if $d\geq (\tau-1)n+2\tau-1$ there exists a connected component of $\Upsilon_{d,n}^t$ such that the polarization of a general element $(X,H)$ is base point free and $H^{n+2}$ is very ample.
    \end{itemize}
These hold unless the considered moduli space is empty.
\end{corollary}




\noindent{\bf Acknowledgments.} 
The author is grateful to Prof. Ljudmila Kamenova for introducing and supporting the project. I also thanks Prof. Christian Schnell for their conversations regarding results in this paper.


\bibliographystyle{plain}
{\small
\bibliography{ref}}

\noindent {\sc Alessandro Pilastro\\
Department of Mathematics\\
Stony Brook University \\
Stony Brook, NY 11794-3651, USA,} \\
\tt alessandro.pilastro@stonybrook.edu\\
\end{document}
\typeout{get arXiv to do 4 passes: Label(s) may have changed. Rerun}
